% ======================================================================
% Beamer presentation - generated by KherveSlide.
% Edit freely; the source is regenerated when you change a slide.
% ======================================================================
\documentclass[aspectratio=169]{beamer}

% --- theme & packages ---
\usetheme{Madrid}
\usepackage{lmodern}
\usepackage{tikz}
\usetikzlibrary{arrows.meta}
\usepackage{colortbl}
\definecolor{ksTblHead}{HTML}{FCE4D6}
\definecolor{ksTblHeadFg}{HTML}{C55A11}
\definecolor{ksTblRule}{HTML}{F4B183}
\definecolor{ksTblCap}{HTML}{808080}
\usepackage[absolute,overlay]{textpos}
\usepackage{graphicx}
\setlength{\TPHorizModule}{1\paperwidth}
\setlength{\TPVertModule}{1\paperheight}
\textblockorigin{0\paperwidth}{0\paperheight}
\setbeamerfont{itemize/enumerate subbody}{size=\relax}
\setbeamerfont{itemize/enumerate subsubbody}{size=\relax}
\setbeamertemplate{footline}{%
\leavevmode\hbox{%
\begin{beamercolorbox}[wd=.333\paperwidth,ht=2.5ex,dp=1.2ex,leftskip=1.5ex]{author in head/foot}\end{beamercolorbox}%
\begin{beamercolorbox}[wd=.334\paperwidth,ht=2.5ex,dp=1.2ex,center]{title in head/foot}\end{beamercolorbox}%
\begin{beamercolorbox}[wd=.333\paperwidth,ht=2.5ex,dp=1.2ex,rightskip=1.5ex]{date in head/foot}\hfill \insertframenumber\,/\,\inserttotalframenumber\end{beamercolorbox}}%
\vskip0pt}
\title{Thermodynamics, lecture 4}
\author{Dr. Teacher}

\begin{document}

% ======================================================================
% Slide 1
% ======================================================================
\begin{frame}[t]

  % text box (free-positioned)
\begin{textblock}{0.84}(0.08,0.16)
{\centering\fontsize{32}{38}\selectfont \textcolor[HTML]{1A1A1A}{\textbf{Thermodynamics, lecture 4}}\par}
\end{textblock}

  % line / arrow (free-positioned)
\definecolor{ksline0}{HTML}{1F4E79}
\begin{textblock}{0.24}(0.38,0.52)
\begin{tikzpicture}
\useasboundingbox (0,0) rectangle (\linewidth,-0\TPVertModule);
\draw[line width=2.5pt,color=ksline0] (0\linewidth,-0\TPVertModule) -- (1\linewidth,-0\TPVertModule);
\end{tikzpicture}
\end{textblock}

  % text box (free-positioned)
\begin{textblock}{0.84}(0.08,0.57)
{\centering\fontsize{17}{20}\selectfont \textcolor[HTML]{555555}{The first law and its applications}\par}
\end{textblock}

  % text box (free-positioned)
\begin{textblock}{0.84}(0.08,0.76)
{\centering\fontsize{13}{16}\selectfont \textcolor[HTML]{555555}{Physics 101  ·  Week 4}\par}
\end{textblock}
\end{frame}

% ======================================================================
% Slide 2 - Learning objectives
% ======================================================================
\begin{frame}[t]
  \frametitle{Learning objectives}

  % text box (free-positioned)
\begin{textblock}{0.88}(0.06,0.22)
{\raggedright\fontsize{17}{20}\selectfont By the end of this lecture you can:\par}
\end{textblock}

  % text box (free-positioned)
\begin{textblock}{0.88}(0.06,0.3)
{\raggedright\fontsize{18}{22}\selectfont \begin{itemize}
\item State the first law of thermodynamics
\item Tell work from heat
\begin{itemize}
\item sign conventions
\item path dependence
\end{itemize}
\item Apply it to ideal-gas processes
\end{itemize}\par}
\end{textblock}
\end{frame}

% ======================================================================
% Slide 3 - The first law
% ======================================================================
\begin{frame}[t]
  \frametitle{The first law}

  % text box (free-positioned)
\begin{textblock}{0.88}(0.06,0.22)
{\begin{definition}\raggedright\fontsize{15}{18}\selectfont The internal energy $U$ of a system is a state function.\end{definition}\par}
\end{textblock}

  % text box (free-positioned)
\begin{textblock}{0.88}(0.06,0.42)
{\begin{theorem}[First law]\raggedright\fontsize{15}{18}\selectfont For any process, the change in internal energy equals the heat added minus the work done by the system.\end{theorem}\par}
\end{textblock}

  % text box (free-positioned)
\begin{textblock}{0.5}(0.25,0.66)
{\centering\fontsize{30}{36}\selectfont \[\Delta U = Q - W\]\par}
\end{textblock}
\end{frame}

% ======================================================================
% Slide 4 - Worked example
% ======================================================================
\begin{frame}[t]
  \frametitle{Worked example}

  % text box (free-positioned)
\begin{textblock}{0.42}(0.04,0.22)
{\begin{block}{Problem}\raggedright\fontsize{14}{17}\selectfont One mole of ideal gas expands isothermally at 300 K from 10 L to 20 L. Find $Q$ and $W$.\end{block}\par}
\end{textblock}

  % text box (free-positioned)
\begin{textblock}{0.46}(0.5,0.22)
{\raggedright\fontsize{15}{18}\selectfont \begin{enumerate}
\item Isothermal: $\Delta U = 0$, so $Q = W$
\item $W = nRT\ln(V_2/V_1)$
\item $W = 8.314 \times 300 \times \ln 2$
\item $W = Q \approx 1.73$ kJ
\end{enumerate}\par}
\end{textblock}

  % text box (free-positioned)
\begin{textblock}{0.42}(0.04,0.6)
{\centering\fontsize{20}{24}\selectfont \[W=\int_{V_1}^{V_2} p\,dV\]\par}
\end{textblock}
\end{frame}

% ======================================================================
% Slide 5 - Typical values
% ======================================================================
\begin{frame}[t]
  \frametitle{Typical values}

  % table (free-positioned)
\begin{textblock}{0.84}(0.08,0.22)
{\definecolor{ksTH3333B3}{HTML}{3333B3}\definecolor{ksTFFFFFFF}{HTML}{FFFFFF}\definecolor{ksTRC8C8C8}{HTML}{C8C8C8}\definecolor{ksTSEEEEF8}{HTML}{EEEEF8}\setlength{\arrayrulewidth}{0.8pt}\arrayrulecolor{ksTRC8C8C8} \fontsize{14}{17}\selectfont \begin{tabular}{|c|c|c|c|}
\hline
  \rowcolor{ksTH3333B3}\textcolor{ksTFFFFFFF}{\textbf{Process}} & \textcolor{ksTFFFFFFF}{\textbf{Constant}} & \textcolor{ksTFFFFFFF}{\textbf{$\Delta U$}} & \textcolor{ksTFFFFFFF}{\textbf{$W$}} \\
\hline
  Isothermal & $T$ & 0 & $nRT\ln(V_2/V_1)$ \\
\hline
  \rowcolor{ksTSEEEEF8}Isobaric & $p$ & $nC_v\Delta T$ & $p\,\Delta V$ \\
\hline
  Isochoric & $V$ & $nC_v\Delta T$ & 0 \\
\hline
  \rowcolor{ksTSEEEEF8}Adiabatic & $Q=0$ & $-W$ & $\frac{p_1V_1-p_2V_2}{\gamma-1}$ \\
\hline
\end{tabular}\\[2pt]{\footnotesize\itshape\textcolor{ksTblCap}{Ideal gas, $n$ moles}}}
\end{textblock}
\end{frame}

% ======================================================================
% Slide 6 - Check your understanding
% ======================================================================
\begin{frame}[t]
  \frametitle{Check your understanding}

  % text box (free-positioned)
\begin{textblock}{0.88}(0.06,0.22)
{\raggedright\fontsize{17}{20}\selectfont Gas is compressed adiabatically. Its temperature…\par}
\end{textblock}

  % shape (free-positioned)
\definecolor{ksfill1}{HTML}{2E7D4F}%
\begin{textblock}{0.25}(0.08,0.36)
\begin{tikzpicture}
\useasboundingbox (0,0) rectangle (\linewidth,-0.22\TPVertModule);
\path[fill=ksfill1, rounded corners=6pt] (0\linewidth,-0\TPVertModule) rectangle (1\linewidth,-0.22\TPVertModule);
\end{tikzpicture}
\end{textblock}

  % text box (free-positioned)
\begin{textblock}{0.25}(0.08,0.3805)
{\centering\fontsize{26}{31}\selectfont \textcolor[HTML]{FFFFFF}{\textbf{A}}\par}
\end{textblock}

  % text box (free-positioned)
\begin{textblock}{0.25}(0.08,0.4869)
{\centering\fontsize{14}{17}\selectfont \textcolor[HTML]{FFFFFF}{rises}\par}
\end{textblock}

  % shape (free-positioned)
\definecolor{ksfill2}{HTML}{1F4E79}%
\begin{textblock}{0.25}(0.37,0.36)
\begin{tikzpicture}
\useasboundingbox (0,0) rectangle (\linewidth,-0.22\TPVertModule);
\path[fill=ksfill2, rounded corners=6pt] (0\linewidth,-0\TPVertModule) rectangle (1\linewidth,-0.22\TPVertModule);
\end{tikzpicture}
\end{textblock}

  % text box (free-positioned)
\begin{textblock}{0.25}(0.37,0.3805)
{\centering\fontsize{26}{31}\selectfont \textcolor[HTML]{FFFFFF}{\textbf{B}}\par}
\end{textblock}

  % text box (free-positioned)
\begin{textblock}{0.25}(0.37,0.4869)
{\centering\fontsize{14}{17}\selectfont \textcolor[HTML]{FFFFFF}{stays the same}\par}
\end{textblock}

  % shape (free-positioned)
\definecolor{ksfill3}{HTML}{B03A3A}%
\begin{textblock}{0.25}(0.66,0.36)
\begin{tikzpicture}
\useasboundingbox (0,0) rectangle (\linewidth,-0.22\TPVertModule);
\path[fill=ksfill3, rounded corners=6pt] (0\linewidth,-0\TPVertModule) rectangle (1\linewidth,-0.22\TPVertModule);
\end{tikzpicture}
\end{textblock}

  % text box (free-positioned)
\begin{textblock}{0.25}(0.66,0.3805)
{\centering\fontsize{26}{31}\selectfont \textcolor[HTML]{FFFFFF}{\textbf{C}}\par}
\end{textblock}

  % text box (free-positioned)
\begin{textblock}{0.25}(0.66,0.4869)
{\centering\fontsize{14}{17}\selectfont \textcolor[HTML]{FFFFFF}{falls}\par}
\end{textblock}

  % text box (free-positioned)
\begin{textblock}{0.88}(0.06,0.68)
{\centering\fontsize{14}{17}\selectfont \textcolor[HTML]{666666}{\textit{Hint: $Q = 0$, $W < 0$.}}\par}
\end{textblock}
\end{frame}

% ======================================================================
% Slide 7 - Summary
% ======================================================================
\begin{frame}[t]
  \frametitle{Summary}

  % text box (free-positioned)
\begin{textblock}{0.88}(0.06,0.24)
{\raggedright\fontsize{19}{23}\selectfont \begin{itemize}
\item $\Delta U = Q - W$ — energy is conserved
\item $U$ depends on the state, $Q$ and $W$ on the path
\item Four ideal-gas processes, four sets of results
\end{itemize}\par}
\end{textblock}

  % text box (free-positioned)
\begin{textblock}{0.88}(0.06,0.7)
{\centering\fontsize{15}{18}\selectfont \textcolor[HTML]{1F4E79}{\textit{Next week: the second law and entropy.}}\par}
\end{textblock}
\end{frame}

\end{document}
